% NR56081C
% Krampfender Patient
% 14.0
% 2013-11-30

\label{m:zerebraler-krampfanfall-bestendend}

\Ca{Vitale Bedrohung!}

\EE{Taktik}: \Speech{Verletzungen vermeiden,
Vitalfunktionen sichern, Abklären ob Status epilepticus
besteht, Klärung des Herganges (Fremdanamnese!),
Differentialdiagnosen abklären}

\begin{itemize}
  \item |TxMassVitMKBes|
  \begin{itemize}
    \item Lagerung: sichern. In der Nachschlafphase stabile
    Seitenlage
    \begin{itemize}
      \item Patient vor Verletzungen schützen, Möbel etc. entfernen oder polstern
      \item NICHT festhalten -- aber sichern, Eigenschutz
      beachten!\footnote{\textbf{Sichern}: Patient darf
      nicht von der Trage oder aus dem Bett fallen, der Kopf
      darf nicht am Boden oder Tischbein anschlagen, etc.}
    \end{itemize}
  \end{itemize}
  \item Besondere Diagnostik:
  \begin{itemize}
    \item Neurocheck inkl. Messung des \EE{Blutzuckers}!
    \item Traumacheck
    \item Nach \E{jedem Krampanfall}: Einschätzungsblock
    (\FullPrettyRef{m:einschaetzungsblock}) inkl.
    Trauma-Untersuchungen
  \end{itemize}
  \item Reizabgeschirmter Transport
  \item Transportentscheidung: Abt.f. Innere Medizin. (evtl. auch
  Intensivüberwachung)
\end{itemize}

\Customized{
\begin{description}
  \item[\CustAsboe] Beim bestehenden zerebralen Krampfanfall ist die Gabe
  eines Benzodiazepins gem. Algorithmus durch NKA vorgesehen.
  \cite{Asb921Memo-AL-2011-000001}.
\end{description}
}